\section{Assumptions, measurability, and nonclaims}
\label{app:measurability}

The measurable representation is an access convention, not a smoothness or
convexity assumption.  The underlying covariate space is the one fixed
nonempty infinite standard Borel space \((\mathcal X,\mathcal A)\) specified
in Appendix~\ref{app:model}; the observation space carries the corresponding
product sigma-field, and the coordinates \(X,T,Y\) and all represented
evaluation maps are measurable.  Countably infinite and uncountable standard
Borel hosts contain the measurable finite blocks needed by every lower
construction.  Finite ambient spaces can have smaller minimax risk, and no
matching lower claim is made for finite ambient spaces.  The E and S arguments
themselves remain valid on every nonempty standard Borel space, but the
headline matching I theorem uses the fixed-infinite field.  Conditional
expectations are versions on this declared probability space.

The residual restriction is marginal: \(\E u^2\in[v_-,v_+]\).
Conditional treatment-residual variance may vary and vanish pointwise.
The upper curvature uses \(\E u^2\), while the lower witnesses in
Appendices~\ref{app:qp}--\ref{app:assembly} retain constant conditional
variance.  All declared pointwise envelopes remain in force.

For any preassigned atomless standard Borel marginal \(\nu\),
Appendix~\ref{app:fixed-nu} and Proposition~\ref{prop:fixed-nu-lower}
give the same-pair lower and matching minimax bounds with constants
independent of \(\nu\).  The public pair may depend on \(\nu,n\) and the
budgets; the rule may know \(\nu\).  The transported saturation witness
retains the original constants.

Proposition~\ref{prop:envelope-bridge} gives the public rule
\(\Pi_{[-3,3]}[P_N(TY)/B_0^2]\), whose expected error is at most
\(1/\sqrt{2n}+4\omega\), where \(\omega=1-v_-/B_0^2\).
This uniform bridge uses no budgets, M1, or \(v_+\), and still requires
the bounded model.  No unbounded heteroskedastic extension is claimed.

For every full-class function, sequential M1 supplies a represented sequence
converging pointwise everywhere inside the fixed \(B_{\mathcal G}\) envelope.
It preserves bounded empirical objectives exactly in the limit and preserves
population \(L_2\) infima by dominated convergence against the
\(B_0\)-bounded truths.
This sequential approximation does not make an arbitrary full-class exact localized supremum measurable.  Accordingly, the model-field certificate
uses outer expectation on each full difference class, while the upper proof uses only the ordinary expectation of measurable countable represented
difference cores, dominated pointwise by the full-class supremum.  For the
finite lower classes the supremum is measurable, so outer expectation and
ordinary expectation are equal.  No equality between a general full-class
localized complexity and its represented-core complexity is asserted.
Countable infima, first-within-tolerance indices, finite arithmetic, and
clipping are Borel.  The branch objectives over rational finite-dimensional
parameters are objective-dense because the relevant bounded convex functions
are continuous and coercive on the reduced domain.  M1 means that these
countable representations and evaluation access are public.  Q1 means that
for each \(n\) and declared budget cell the pair is fixed before data and any
latent hard label, after which a pair-specific Borel rule is chosen; neither
realized-data classes nor one pair fixed across all \(n\) are required.

\subsection{Why the two strengthened fields are necessary}

The finite-space exclusion is substantive.  On a singleton covariate space,
pair the observations across the two independent halves and put
\(D_i=T_{1i}-T_{2i}\), \(E_i=Y_{1i}-Y_{2i}\), and
\(S_n=\sum_iD_i^2\).  The clipped ratio
\[
 \widehat\beta=\Pi_{[-3,3]}\left[
 \frac{\sum_iD_iE_i}{S_n}\mathbf 1\{S_n>0\}\right]
\]
is Borel and satisfies, uniformly over every compatible law and public pair,
\[
 \E|\widehat\beta-\beta_0|
 \le B_0\sqrt{\frac{2}{nv_-}}
 +6\exp\!\left\{-\frac{nv_-^2}{8B_0^4}\right\}.
 \tag{I.1}
\]
Indeed, differencing removes both constant nuisances, the residual numerator
is conditionally centered with conditional second moment at most
\(2B_0^2S_n\), and Hoeffding gives
\(\Pr(S_n<nv_-)\le\exp\{-nv_-^2/(8B_0^4)\}\).  Yet the legal zero singleton
pair in the cell \((a,s,b,t)=(0,1/2,0,1/2)\) has
\(X=1/16+1/4=5/16\).  Thus a positive fixed multiple of \(q+X\) cannot be a
uniform lower bound on that finite ambient space.

The sequential premise is also substantive.  On \([0,1]\) with Lebesgue
law, let \(R\) consist of \(g_B=\mathbf 1_{[0,1]\setminus B}\), where \(B\)
ranges over finite unions of rational-endpoint intervals of measure at most
\(1/2\), and let \(\cG=R\cup\{0\}\).  Every finite-coordinate neighborhood
of zero meets \(R\), but
\[
 \inf_{r\in R}\norm r_2^2\ge\frac12
 \quad\text{whereas}\quad
 \inf_{g\in\cG}\norm g_2^2=0.
 \tag{I.2}
\]
No sequence in \(R\) converges pointwise to zero: dominated convergence would
force \(P(B_k)\to1\), contradicting \(P(B_k)\le1/2\).  Hence a closure phrase
without the explicit per-element sequence does not supply the comparators
used in Appendix~\ref{app:upper-initial} or the calibration suprema.

The baseline minimax risk permits known pre-data upper certificates for the four budgets, an arbitrary Borel decision rule,
expected absolute loss, \(2n\) iid observations, and constants depending on
fixed \(\rho\), including fixed \(B_{\mathcal G}>0\) and
\(c_{\rm rad}>0\).  It does not establish
one pair-independent estimator, finite-query optimization,
a polynomial-time algorithm, a high-probability confidence interval,
empirical superiority, or universal constants over \(\rho\).  In particular,
constants may degenerate as \(B_{\mathcal G}\downarrow0\) or
\(c_{\rm rad}\downarrow0\), where the public class field collapses and the
headline nuisance lower is not asserted.  The finite lower
instances satisfy the same M1 interface and do not use inaccessible latent
labels as learner inputs.  The target is a statistical PLM coefficient; a
causal or potential-outcomes interpretation would require additional
identification assumptions and is not claimed.
Theorem~\ref{thm:approximation-main} and
Appendix~\ref{app:stochastic-adaptation} strengthen the upper quantifiers:
for each fixed $(\Gamma,\rho,n)$ one rule attains $q+W$ simultaneously
over all four unknown budgets. Empirical local Rademacher fixed points
on an independent covariate block supply the stochastic inputs to the
double-calibrated rule of Appendix~\ref{app:approximation-adaptation}.
A positive-radius margin transfers the core's linear certificate to
the original full outer complexity; the public classes are unchanged.
Separately,
Proposition~\ref{prop:histogram-adaptation} and Appendix~\ref{app:budget-adaptation}
give an all-budget-independent rule and the matching restricted outer-pair
rate \(q+(a+s^2)(b+t^2)\) for full interval boxes on possibly different
partitions around arbitrary public center functions. Theorem~\ref{thm:full-box-minimax}
uses one pair depending on \(\rho,n,s,t\), uniform over \(a,b\).
Proposition~\ref{prop:hilbert-main} and Appendix~\ref{app:hilbert-adaptation}
give the smaller matching outer-pair rate \(q+ab\) for full affine Hilbert
balls, again without budget inputs; zero-radius balls witness its lower.
Both statements impose geometry beyond M1, and neither asserts a lower
for every fixed pair. Finite implementations require explicitly evaluable
centers and the stated partition/width or kernel interface.


Antecedents include partially linear and orthogonal estimation
\citep{robinson1988root,chernozhukov2018doubledebiased,foster2023orthogonal},
structure-agnostic product bounds \citep{balakrishnan2023fundamentalstructureagnostic},
two-mixture testing \citep{robins2009semiparametric}, and Gu's four-budget
question \citep{gu2025structureagnosticplm}.
